\documentclass[10pt, a5paper]{scrartcl}

\usepackage{pylabUkr}

\begin{document}

\myheader{Варіант 5}

\begin{enumerate}
%\item Ввести с клавиатуры целое число (переменная \verb|n1|), вещественное число (переменная \verb|n2|) и слово (переменная \verb|w|).
%\begin{enumerate}
%\item Вывести значения переменных в одной строке, разделяя пробелом.
%\item Вывести значения переменных одно под другим.
%\item Вывести на экран значения переменных, сопровождая вывод комментариями, например:
%\begin{verbatim}
%Вы ввели три значения.
%Первое из них - ...
%Второе - ...
%А третье - ...
%\end{verbatim}
%\end{enumerate}

\item \textbf{Улюблений фільм.} Складіть програму, яка  запитує у користувача назву його улюбленого фільму

(приклад: \verb|Ваш улюблений фільм:|),

а потім, після введення даних, друкує на екрані повідомлення

(приклад: \verb|Мій улюблений фільм - ...|).

\item \textbf{Пиріжок.} Пиріжок у місцевій крамниці коштує $a$ доларів $b$ центів. Скільки доларів та центів доведеться заплатити за $n$ пиріжків?

{\itshape Параметри $a$, $b$ та $n$ вводяться з клавіатури.}

\item \textbf{Упертий Біллі.} Юний Білл Гейтс не любив уроки граматики в школі і якось вирішив втекти з уроку через вікно. Кватирка має форму прямокутника розміром $ m \times n $ сантиметрів. Визначити, чи зможе голова Білла Гейтса пролізти у кватирку, якщо голова є ідеальною кулею діаметром $d$ сантиметрів.

{\itshape Параметри $m$, $n$ та $d$ вводяться з клавіатури.}

\item \textbf{Попередній день.} Дата деякого дня визначається трьома натуральними числами: $d$ (день), $m$ (місяць), $y$ (рік). Визначте дату попереднього дня. Рік не є високосним.

{\itshape Параметри $d$, $m$ та $y$ вводяться з клавіатури.}

% \item \textbf{Перестановка цифр.} Дано четырёхзначное натуральное число. Две его последние цифры были переставлены в начало числа. Вывести полученное число.

% \item \textbf{Формула.} Даны некоторые действительные числа $x$ и $y$. Вычислить значение выражения
% $$
% z=\dfrac{\sqrt{\mathstrut 1+\ln \left(1+x^2\right)}+\sqrt{\mathstrut 1+\dfrac{1}{1+|y|}}}{2x\,(1+y^2)}
% $$
% с точностью до сотых. Проверка данных на корректность не требуется. 
\end{enumerate}

\end{document}

\cleardoublepage

\myheader{Вариант 6}

\begin{enumerate}
%\item Ввести с клавиатуры целое число (переменная \verb|n1|), вещественное число (переменная \verb|n2|) и слово (переменная \verb|w|).
%\begin{enumerate}
%\item Вывести значения переменных в одной строке, разделяя пробелом.
%\item Вывести значения переменных одно под другим.
%\item Вывести на экран значения переменных, сопровождая вывод комментариями, например:
%\begin{verbatim}
%Вы ввели три значения.
%Первое из них - ...
%Второе - ...
%А третье - ...
%\end{verbatim}
%\end{enumerate}

\item \textbf{Университет.} Составьте программу, которая  запрашивает у пользователя название вуза, в котором он учится

(пример: \verb|В каком вузе Вы учитесь?|),

а затем, после ввода данных, печатает на экране сообщение

(пример: \verb|Я учусь в ...|).

\item \textbf{Пробежка.} Даны размеры (в метрах) футбольного поля, имеющего форму прямоугольника. Какое расстояние пробежит спортсмен, <<намотавший>> $n$ кругов вокруг этого поля? Число $n$ задаётся с клавиатуры.

{\itshape Старайтесь сопровождать ввод и вывод текстовыми сообщениями.}

\item \textbf{Задом наперёд.} Дано трёхзначное натуральное число. Вывести на экран число, получающееся при прочтении исходного числа справа налево.

\item \textbf{Формула.} Даны некоторые действительные числа $a$ и $b$. Вычислить значение выражения
$$
z=\sin(a^2+b^2)\cdot\dfrac{\dfrac{a}{a^2+25}+b}{\sqrt{\mathstrut b}+\dfrac{a+b}{2}}
$$
с точностью до сотых. Проверка данных на корректность не требуется. 
\end{enumerate}



\cleardoublepage



\myheader{Вариант 7}

\begin{enumerate}
%\item Ввести с клавиатуры целое число (переменная \verb|n1|), вещественное число (переменная \verb|n2|) и слово (переменная \verb|w|).
%\begin{enumerate}
%\item Вывести значения переменных в одной строке, разделяя пробелом.
%\item Вывести значения переменных одно под другим.
%\item Вывести на экран значения переменных, сопровождая вывод комментариями, например:
%\begin{verbatim}
%Вы ввели три значения.
%Первое из них - ...
%Второе - ...
%А третье - ...
%\end{verbatim}
%\end{enumerate}

\item \textbf{Приветствие.} Составьте программу, которая  запрашивает у пользователя имя и фамилию 

(пример: \verb|Как Вас зовут?|),

а затем, после ввода данных, печатает на экране какое-либо приветствие

(пример: \verb|Привет, ... ...!|).

\item \textbf{Градусы и радианы.} Дано некоторое число, обозначающее градусную меру угла. Перевести это значение в радианы.

{\itshape Старайтесь сопровождать ввод и вывод текстовыми сообщениями.}

\item \textbf{Перестановка цифр.} Дано трёхзначное натуральное число. Вывести на экран число, получающееся при перестановке второй и третьей цифр исходного числа.

\item \textbf{Формула.} Даны некоторые действительные числа $x$ и $y$. Вычислить значение выражения
$$
z=\ln \left(2-\sin x^2\right)\cdot\dfrac{x\,(1+x)}{|y|+\dfrac{1}{\sqrt{\mathstrut \mathstrut x^2+10}}}
$$
с точностью до сотых. Проверка данных на корректность не требуется. 
\end{enumerate}



\cleardoublepage



\myheader{Вариант 8}

\begin{enumerate}
%\item Ввести с клавиатуры целое число (переменная \verb|n1|), вещественное число (переменная \verb|n2|) и слово (переменная \verb|w|).
%\begin{enumerate}
%\item Вывести значения переменных в одной строке, разделяя пробелом.
%\item Вывести значения переменных одно под другим.
%\item Вывести на экран значения переменных, сопровождая вывод комментариями, например:
%\begin{verbatim}
%Вы ввели три значения.
%Первое из них - ...
%Второе - ...
%А третье - ...
%\end{verbatim}
%\end{enumerate}

\item \textbf{Любимый герой.} Составьте программу, которая запрашивает у~пользователя имя его любимого мультипликационного героя

(пример: \verb|Ваш любимый мультяшный герой:|),

а затем, после ввода данных, печатает на экране сообщение

(пример: \verb|Мой любимый герой - ...|).

\item \textbf{Градусы Фаренгейта.} Известно значение температуры по шкале Цельсия. Найти соответствующее
значение температуры по шкале Фаренгейта\footnote{Для пересчёта по шкале Фаренгейта необходимо исходное значение температуры умножить на~1.8 и к результату прибавить 32.}.

{\itshape Старайтесь сопровождать ввод и вывод текстовыми сообщениями.}

\item \textbf{Сумма двух цифр.} Дано четырёхзначное натуральное число. Найти сумму двух его последних цифр.

\item \textbf{Формула.} Даны некоторые действительные числа $a$ и $b$. Вычислить значение выражения
$$
z=\cos(a^2+b^2)\cdot\dfrac{|a+b|}{2(b-2a)(a^2+b^2)}
$$
с точностью до сотых. Проверка данных на корректность не требуется. 
\end{enumerate}

\cleardoublepage

\myheader{Вариант 9}

\begin{enumerate}
%\item Ввести с клавиатуры целое число (переменная \verb|n1|), вещественное число (переменная \verb|n2|) и слово (переменная \verb|w|).
%\begin{enumerate}
%\item Вывести значения переменных в одной строке, разделяя пробелом.
%\item Вывести значения переменных одно под другим.
%\item Вывести на экран значения переменных, сопровождая вывод комментариями, например:
%\begin{verbatim}
%Вы ввели три значения.
%Первое из них - ...
%Второе - ...
%А третье - ...
%\end{verbatim}
%\end{enumerate}

\item \textbf{Любимый цвет.} Составьте программу, которая  запрашивает у~пользователя его любимый цвет 

(пример: \verb|Назовите Ваш любимый цвет:|),

а затем, после ввода данных, печатает на экране сообщение

(пример: \verb|Мой любимый цвет - ...|).

\item \textbf{Плотность населения.} Известны количество жителей в городе и площадь его территории (в км$^2$). Определить плотность населения в этом городе\footnote[1]{По состоянию на июль 2018 года городом с наибольшей плотностью населения является Манила "--- столица государства Филиппины (43\,079 человек на км$^2$).}.

{\itshape Старайтесь сопровождать ввод и вывод текстовыми сообщениями.}

\item \textbf{Произведение цифр.} Дано трёхзначное натуральное число. Найти произведение его цифр.

\item \textbf{Формула.} Даны некоторые действительные числа $a$ и $b$. Вычислить значение выражения
$$
z=\cos(a^2+b^2)\cdot\dfrac{|a+b|}{2(b-2a)(a^2+b^2)}
$$
с точностью до сотых. Проверка данных на корректность не требуется. 
\end{enumerate}

\cleardoublepage

\myheader{Вариант 10}

\begin{enumerate}
%\item Ввести с клавиатуры целое число (переменная \verb|n1|), вещественное число (переменная \verb|n2|) и слово (переменная \verb|w|).
%\begin{enumerate}
%\item Вывести значения переменных в одной строке, разделяя пробелом.
%\item Вывести значения переменных одно под другим.
%\item Вывести на экран значения переменных, сопровождая вывод комментариями, например:
%\begin{verbatim}
%Вы ввели три значения.
%Первое из них - ...
%Второе - ...
%А третье - ...
%\end{verbatim}
%\end{enumerate}

\item \textbf{Любимый мультфильм.} Составьте программу, которая  запрашивает у~пользователя название любимого мультфильма 

(пример: \verb|Ваш любимый мультфильм:|),

а затем, после ввода данных, печатает на экране сообщение

(пример: \verb|Мой любимый мультфильм - ...|).

\item \textbf{Квадрат и точки.} Какое наибольшее количество точек с целочисленными координатами можно на листке в клеточку накрыть квадратом со стороной $n$ клеток? Число $n$ задаётся с клавиатуры.

{\itshape Старайтесь сопровождать ввод и вывод текстовыми сообщениями.}

\item \textbf{Номер билета.} Дано шестизначное натуральное число "--- номер проездного билета. Найти сумму цифр в <<левой>> и <<правой>> частях номера.

\item \textbf{Формула.} Даны некоторые действительные числа $a$ и $b$. Вычислить значение выражения
$$
z=\sin(a^2-b^2)\cdot\dfrac{|a+b|}{2(b-2a)(a^2-b^2)}
$$
с точностью до сотых. Проверка данных на корректность не требуется. 
\end{enumerate}

\cleardoublepage



\myheader{Вариант 11}

\begin{enumerate}
%\item Ввести с клавиатуры целое число (переменная \verb|n1|), вещественное число (переменная \verb|n2|) и слово (переменная \verb|w|).
%\begin{enumerate}
%\item Вывести значения переменных в одной строке, разделяя пробелом.
%\item Вывести значения переменных одно под другим.
%\item Вывести на экран значения переменных, сопровождая вывод комментариями, например:
%\begin{verbatim}
%Вы ввели три значения.
%Первое из них - ...
%Второе - ...
%А третье - ...
%\end{verbatim}
%\end{enumerate}

\item \textbf{Любимый киногерой.} Составьте программу, которая запрашивает у~пользователя имя его любимого киногероя

(пример: \verb|Ваш любимый киногерой:|),

а затем, после ввода данных, печатает на экране сообщение

(пример: \verb|Мой любимый киногерой - ...|).

\item \textbf{Углы треугольника.} Даны катеты прямоугольного треугольника. Найти все его углы (в радианах).

{\itshape Старайтесь сопровождать ввод и вывод текстовыми сообщениями.}

\item \textbf{Перестановка цифр.} Дано трёхзначное натуральное число. Вывести на экран число, получающееся при перестановке первой и второй  цифр исходного числа.

\item \textbf{Формула.} Даны некоторые действительные числа $a$ и $b$. Вычислить значение выражения
$$
z=\cos(a^2+b^2)\cdot\dfrac{|a+b|}{2(b-2a)(a^2+b^2)}
$$
с точностью до сотых. Проверка данных на корректность не требуется. 
\end{enumerate}

\cleardoublepage

\myheader{Вариант 12}

\begin{enumerate}
%\item Ввести с клавиатуры целое число (переменная \verb|n1|), вещественное число (переменная \verb|n2|) и слово (переменная \verb|w|).
%\begin{enumerate}
%\item Вывести значения переменных в одной строке, разделяя пробелом.
%\item Вывести значения переменных одно под другим.
%\item Вывести на экран значения переменных, сопровождая вывод комментариями, например:
%\begin{verbatim}
%Вы ввели три значения.
%Первое из них - ...
%Второе - ...
%А третье - ...
%\end{verbatim}
%\end{enumerate}

\item \textbf{День недели.} Составьте программу, которая  запрашивает у пользователя день недели

(пример: \verb|Какой сегодня день недели?|),

а затем, после ввода данных, печатает на экране сообщение

(пример: \verb|Ура, сегодня ...!|).

\item \textbf{Треугольник.} Даны катеты прямоугольного треугольника. Найти его углы (в градусах).

{\itshape Старайтесь сопровождать ввод и вывод текстовыми сообщениями.}

\item \textbf{Перестановка цифр.} Дано четырёхзначное натуральное число. Поменять местами его первую и последнюю цифры. Вывести полученное число.

\item \textbf{Формула.} Даны некоторые действительные числа $x$ и $y$. Вычислить значение выражения
$$
z=\dfrac{\sqrt{\mathstrut 1+\ln \left(1+x^2\right)}+\sqrt{\mathstrut 1+\dfrac{1}{1+|y|}}}{2x\,(1+y^2)}
$$
с точностью до сотых. Проверка данных на корректность не требуется. 
\end{enumerate}

\cleardoublepage

\myheader{Вариант 13}

\begin{enumerate}
%\item Ввести с клавиатуры целое число (переменная \verb|n1|), вещественное число (переменная \verb|n2|) и слово (переменная \verb|w|).
%\begin{enumerate}
%\item Вывести значения переменных в одной строке, разделяя пробелом.
%\item Вывести значения переменных одно под другим.
%\item Вывести на экран значения переменных, сопровождая вывод комментариями, например:
%\begin{verbatim}
%Вы ввели три значения.
%Первое из них - ...
%Второе - ...
%А третье - ...
%\end{verbatim}
%\end{enumerate}

\item \textbf{Любимая группа.} Составьте программу, которая  запрашивает у пользователя название его любимой музыкальной группы

(пример: \verb|Ваша любимая музыкальная группа:|),

а затем, после ввода данных, печатает на экране сообщение

(пример: \verb|Моя любимая группа - ...|).

\item \textbf{Утренняя пробежка.} Даны размеры (в метрах) футбольного поля, имеющего форму прямоугольника. Сколько кругов нужно <<намотать>> вокруг этого поля, чтобы пробежать суммарно $n$ км? Число $n$ задаётся с клавиатуры.

{\itshape Старайтесь сопровождать ввод и вывод текстовыми сообщениями.}

\item \textbf{Задом наперёд.} Дано трёхзначное натуральное число. Вывести на экран число, получающееся при прочтении исходного числа справа налево.

\item \textbf{Формула.} Даны некоторые действительные числа $a$ и $b$. Вычислить значение выражения
$$
z=\cos(a^2-b^2)\cdot\dfrac{\dfrac{a}{a^2+25}+b}{\sqrt{\mathstrut b}+\dfrac{a+b}{3}}
$$
с точностью до сотых. Проверка данных на корректность не требуется. 
\end{enumerate}

\cleardoublepage

\myheader{Вариант 14}

\begin{enumerate}
%\item Ввести с клавиатуры целое число (переменная \verb|n1|), вещественное число (переменная \verb|n2|) и слово (переменная \verb|w|).
%\begin{enumerate}
%\item Вывести значения переменных в одной строке, разделяя пробелом.
%\item Вывести значения переменных одно под другим.
%\item Вывести на экран значения переменных, сопровождая вывод комментариями, например:
%\begin{verbatim}
%Вы ввели три значения.
%Первое из них - ...
%Второе - ...
%А третье - ...
%\end{verbatim}
%\end{enumerate}

\item \textbf{Любимый город.} Составьте программу, которая  запрашивает у пользователя название любимого города

(пример: \verb|Какой Ваш любимый город?|),

а затем, после ввода данных, печатает на экране какое-либо приветствие

(пример: \verb|Мой любимый город - ...|).

\item \textbf{Бег по кругу.} Леди Гамильтон "--- очень спортивная женщина, несмотря на почтенный возраст. Её любимое занятие по утрам "--- пробежка вокруг знаменитого фонтана в парке Кентербери вместе с верным псом по кличке Линкольн. Какое расстояние (в метрах) обычно пробегает по утрам миссис Гамильтон с Линкольном, если радиус фонтана $r$ метров и они делают $n$ кругов?

{\itshape Старайтесь сопровождать ввод и вывод текстовыми сообщениями.}

\item \textbf{Сумма квадратов цифр.} Дано трёхзначное натуральное число. Вывести на экран сумму квадратов его цифр.

\item \textbf{Формула.} Даны некоторые действительные числа $x$ и $y$. Вычислить значение выражения
$$
z=\ln \left(2-\sin x^2\right)\cdot\dfrac{x\,(1+x)}{|y|+\dfrac{1}{\sqrt{\mathstrut \mathstrut x^2+10}}}
$$
с точностью до сотых. Проверка данных на корректность не требуется. 
\end{enumerate}

\cleardoublepage

\myheader{Вариант 15}

\begin{enumerate}
%\item Ввести с клавиатуры целое число (переменная \verb|n1|), вещественное число (переменная \verb|n2|) и слово (переменная \verb|w|).
%\begin{enumerate}
%\item Вывести значения переменных в одной строке, разделяя пробелом.
%\item Вывести значения переменных одно под другим.
%\item Вывести на экран значения переменных, сопровождая вывод комментариями, например:
%\begin{verbatim}
%Вы ввели три значения.
%Первое из них - ...
%Второе - ...
%А третье - ...
%\end{verbatim}
%\end{enumerate}

\item \textbf{Приветствие.} Составьте программу, которая запрашивает у пользователя его имя и год рождения, а затем выводит сообщение

\verb|Здравствуйте, ...! Ваш год рождения - ...|

\item \textbf{Брутальная клумба.} Найти площадь клумбы, если она представляет собой идеальную окружность, а на её ограждение ушло $n$ метров бетонного забора. Высоту забора не учитываем. Число $n$ вводится с клавиатуры.

{\itshape Старайтесь сопровождать ввод и вывод текстовыми сообщениями.}

\item \textbf{Цифры числа.} Дано четырёхзначное натуральное число. Переставить последнюю цифру в~начало числа. Вывести результат на экран.

\item \textbf{Формула.} Даны некоторые действительные числа $a$ и $b$. Вычислить значение выражения
$$
z=\cos(a^2+b^2)\cdot\dfrac{|a+b|}{2(b-2a)(a^2+b^2)}
$$
с точностью до сотых. Проверка данных на корректность не требуется. 
\end{enumerate}

\cleardoublepage

\myheader{Вариант 16}

\begin{enumerate}
%\item Ввести с клавиатуры целое число (переменная \verb|n1|), вещественное число (переменная \verb|n2|) и слово (переменная \verb|w|).
%\begin{enumerate}
%\item Вывести значения переменных в одной строке, разделяя пробелом.
%\item Вывести значения переменных одно под другим.
%\item Вывести на экран значения переменных, сопровождая вывод комментариями, например:
%\begin{verbatim}
%Вы ввели три значения.
%Первое из них - ...
%Второе - ...
%А третье - ...
%\end{verbatim}
%\end{enumerate}

\item \textbf{Любимый автор.} Составьте программу, которая  запрашивает у~пользователя фамилию его любимого автора книг

(пример: \verb|Ваш любимый автор:|),

а затем, после ввода данных, печатает на экране сообщение

(пример: \verb|Мой любимый автор - ...|).

\item \textbf{Сумма чисел.} Найти сумму всех целых чисел на промежутке $[a; b]$. Числа $a$ и $b$ "--- целые, вводятся с клавиатуры.

{\itshape Старайтесь сопровождать ввод и вывод текстовыми сообщениями.}

\item \textbf{Сумма двух цифр.} Дано четырёхзначное натуральное число. Найти сумму двух его последних цифр.

\item \textbf{Формула.} Даны некоторые действительные числа $a$ и $b$. Вычислить значение выражения
$$
z=\sin(a^2+b^2)\cdot\dfrac{\dfrac{a}{a^2+25}+b}{\sqrt{\mathstrut b}+\dfrac{a+b}{2}}
$$
с точностью до сотых. Проверка данных на корректность не требуется. 
\end{enumerate}

\cleardoublepage

\myheader{Вариант 17}

\begin{enumerate}
%\item Ввести с клавиатуры целое число (переменная \verb|n1|), вещественное число (переменная \verb|n2|) и слово (переменная \verb|w|).
%\begin{enumerate}
%\item Вывести значения переменных в одной строке, разделяя пробелом.
%\item Вывести значения переменных одно под другим.
%\item Вывести на экран значения переменных, сопровождая вывод комментариями, например:
%\begin{verbatim}
%Вы ввели три значения.
%Первое из них - ...
%Второе - ...
%А третье - ...
%\end{verbatim}
%\end{enumerate}

\item \textbf{Любимый город.} Составьте программу, которая  запрашивает у пользователя название любимого города

(пример: \verb|Какой Ваш любимый город?|),

а затем, после ввода данных, печатает на экране какое-либо приветствие

(пример: \verb|Мой любимый город - ...|).

\item \textbf{Куб.} Дана длина ребра куба. Найти объём куба и площадь его боковой поверхности.

{\itshape Старайтесь сопровождать ввод и вывод текстовыми сообщениями.}

\item \textbf{Перестановка цифр.} Дано четырёхзначное натуральное число. Поменять местами его первую и последнюю цифры. Вывести полученное число.

\item \textbf{Формула.} Даны некоторые действительные числа $a$ и $b$. Вычислить значение выражения
$$
z=\cos(a^2+b^2)\cdot\dfrac{|a+b|}{2(b-2a)(a^2+b^2)}
$$
с точностью до сотых. Проверка данных на корректность не требуется.  
\end{enumerate}

\cleardoublepage

\myheader{Бонусные задачи}

\begin{enumerate}
\item Даны три числа. Сколько среди них нечётных?

\item Дано произвольное натуральное число. Найти количество цифр в нём.

\item Даны два числа $a$ и $b$. Как поменять их значения местами?

Подумайте, почему не работает <<наивное решение>>: 
\begin{verbatim}
a=b
b=a
\end{verbatim}

\item Попробуйте поменять местами значения $a$ и $b$ (см. предыдущую задачу), не используя дополнительную переменную.

\item Дано целое число $k$ ($1\leqslant k \leqslant 365$). Каким (по счёту) днём недели является $k$-й день года, если 1 января этого года: а) понедельник; б) вторник?

\item Где-то в глубинах космоса затерялась странная планета Кошиньютон. Правят ею лагранжианы "--- гигантские голубые существа с~красными глазами, обладающие сверхвысоким уровнем развития. Чтобы другие цивилизации не надумали вступить с~ними в контакт, они используют весьма запутанный алфавит, содержащий $1\,234\,567\,890^{2}$ символов.
  
Один из местных жителей был отправлен на Землю со шпионской миссией, однако возникла проблема с кодировкой символов лагранжианского алфавита на земном компьютере.

Какое наименьшее число бит понадобится для кодирования всех символов?

\cleardoublepage

%\myheader{Бонусные задачи}
%
%\begin{enumerate}
%\setcounter{enumi}{6}

\item Дано число $n$. С начала суток прошло $n$ минут. Определите, сколько часов и минут будут показывать электронные часы в этот момент. Программа должна вывести два числа: количество часов (от 0 до 23) и количество минут (от 0 до 59). Учтите, что число $n$ может быть больше, чем количество минут в сутках.

\item Даны значения двух моментов времени, принадлежащих одним и тем же суткам: часы, минуты и секунды для каждого из моментов времени. Известно, что второй момент времени наступил не раньше первого. Определите, сколько секунд прошло между двумя моментами времени.

Программа на вход получает три целых числа ($h1$, $m1$, $s1$), задающие первый момент времени и три целых числа ($h2$, $m2$, $s2$), задающих второй момент времени. Выведите число секунд между этими моментами времени.

\item В зоопарке $n$ клеток выстроены в ряд. Кроме прочих обитателей, в нём живут две мартышки, Слава и Юра. Слава и Юра всегда были большими друзьями и сидели в соседних клетках, но теперь они поссорились и больше не хотят видеть друг друга. Смотритель уже собрался переселить их в соответствии с их желанием, однако возникла проблема. Слава и Юра "--- очень образованные мартышки, и они непременно хотят знать, сколько всего существует способов расселить их так, чтобы их клетки не были соседними, и уж конечно их клетки должны быть различными. Можно считать, что все $n$ клеток доступны, остальные обитатели зоопарка готовы переехать куда угодно.

Смотритель попытался было сосчитать сам, но сбился со счёта где-то в районе бегемотов. И совершенно определённо, без вашей помощи ему не справиться!

\item Предприниматель, имея стартовый капитал $k$ евро, занялся торговлей, которая ежемесячно увеличивает капитал на $p$\,\%. Через сколько лет он накопит сумму $S$, достаточную для покупки собственного магазина?

%\item Функция $y=\sin x$ на отрезке $x\in\left[0; \dfrac{\pi}{2}\right]$ хорошо аппроксимируется выражением $x-\dfrac{x^3}{6}+\dfrac{x^5}{120}$\footnote[1]{Ответ на вопрос, откуда взялась эта формула, вы скоро узнаете в курсе математического анализа.}. Для заданного значения аргумента $x$ вычислить $y$ по этой формуле и сравнить с точным значением, вычисленным с помощью стандартной функции \verb|math.sin(x)|.


\end{enumerate}

\end{document}

%%%%

\myheader{Задачи}

\begin{enumerate}
\item Дано число $n$. С начала суток прошло $n$ минут. Определите, сколько часов и минут будут показывать электронные часы в этот момент. Программа должна вывести два числа: количество часов (от 0 до 23) и количество минут (от 0 до 59). Учтите, что число $n$ может быть больше, чем количество минут в сутках.

\item Дано три числа. Сколько среди них нечётных?

\item Дано трёхзначное натуральное число. Найти сумму его цифр.

\item Дана длина ребра куба. Найти объём куба и площадь его боковой поверхности.



\item Дано трёхзначное натуральное число. В нём зачеркнули первую цифру и~приписали её в конце. Вывести на экран полученное число.

\item Найти площадь кольца по заданным внешнему и внутреннему радиусам.



\item Дано четырёхзначное натуральное число. Найти сумму двух его первых цифр.

\item Дано четырёхзначное натуральное число. Переставить первую цифру в~конец числа. Вывести результат на экран.

\item Даны катеты прямоугольного треугольника. Найти его периметр.



\item Дано четырёхзначное натуральное число. Две его последние цифры были переставлены в начало числа. Вывести полученное число.

\item Даны стороны прямоугольника. Найти его периметр и длину диагонали.

\item Дано трёхзначное натуральное число. Вывести на экран число, получающееся при прочтении исходного числа справа налево.

\end{enumerate}

\myheader{Задачи}

\begin{enumerate}
\setcounter{enumi}{12}

\item Дано некоторое число, обозначающее градусную меру угла. Перевести это значение в радианы.



\item Дано трёхзначное натуральное число. Вывести на экран число, получающееся при перестановке второй и третьей цифр исходного числа.

\item Даны некоторые действительные числа $x$ и $y$. Вычислить значение выражения
$$
z=\ln \left(2-\sin x^2\right)\cdot\dfrac{x\,(1+x)}{|y|+\dfrac{1}{\sqrt{\mathstrut \mathstrut x^2+10}}}
$$
с точностью до сотых. Проверка данных на корректность не требуется. 

\item Известно значение температуры по шкале Цельсия. Найти соответствующее
значение температуры по шкале Фаренгейта\footnote{Для пересчёта по шкале Фаренгейта необходимо исходное значение температуры умножить на~1.8 и к результату прибавить 32.}.



\item Дано четырёхзначное натуральное число. Найти сумму двух его последних цифр.

\item Даны некоторые действительные числа $a$ и $b$. Вычислить значение выражения
$$
z=\cos(a^2+b^2)\cdot\dfrac{|a+b|}{2(b-2a)(a^2+b^2)}
$$
с точностью до сотых. Проверка данных на корректность не требуется. 

\item Известны количество жителей в государстве и площадь его территории. Определить плотность населения в этом государстве.

\end{enumerate}

\myheader{Задачи}

\begin{enumerate}
\setcounter{enumi}{19}

\item Дано трёхзначное натуральное число. Вывести на экран произведение его цифр.

\item Пирожок в столовой стоит $a$ долларов $b$ центов. Определите, сколько долларов и центов нужно заплатить за $n$ пирожков. Программа получает на вход три числа: $a$, $b$, $n$ и должна вывести два числа: стоимость покупки в долларах и центах.

\item Даны значения двух моментов времени, принадлежащих одним и тем же суткам: часы, минуты и секунды для каждого из моментов времени. Известно, что второй момент времени наступил не раньше первого. Определите, сколько секунд прошло между двумя моментами времени.
Программа на вход получает три целых числа: часы, минуты, секунды, задающие первый момент времени и три целых числа, задающих второй момент времени.
Выведите число секунд между этими моментами времени.

\item Для настольной игры используются карточки с номерами от 1 до $N$. Одна карточка потерялась. Найдите её, зная номера оставшихся карточек.
Дано число $N$, далее $(N-1)$ номер оставшихся карточек. Программа должна вывести номер потерянной карточки.

\item Вычислить значение выражения 
$$
e^{x-2}+|\sin x|-x^4 \cos\dfrac{1}{x}
$$
при $x=3.6$. \textit{(Ответ: $-156.1276$.)}

\item Дан прямоугольник размером $647 \times 170$. Сколько квадратов со стороной 30 можно вырезать из него?

\end{enumerate}



\end{document}